\documentclass[11pt]{article}
\usepackage[margin=1in]{geometry}
\usepackage{amsmath,amssymb,amsthm}
\usepackage{graphicx}
\usepackage[hidelinks]{hyperref}
\setlength{\emergencystretch}{2em}

\newtheorem{theorem}{Theorem}
\newtheorem{corollary}{Corollary}
\newtheorem{remark}{Remark}

\newcommand{\F}{\mathbb F}
\newcommand{\C}{\mathbb C}
\newcommand{\dist}{\operatorname{dist}}

\title{BCP and Countable Capture Are Exact Opposites}
\author{Run \texttt{fa\_banach\_001}, lane 19}
\date{June 17, 2026}

\begin{document}
\maketitle

\section*{Source Problem}

The source paper is Jinghao Huang, Karimbergen Kudaybergenov, and Rui Liu,
\emph{The ball-covering property of von Neumann algebras and noncommutative
symmetric spaces}, arXiv:2311.14257.  After proving that a commutative
unital \(C^*\)-algebra \(C(K)\) has BCP/UBCP exactly when \(K\) has a
countable \(\pi\)-basis, it asks for a noncommutative version:

\begin{center}
\includegraphics[width=0.92\textwidth]{figures/open_problem_crop.png}
\end{center}

\section*{Definitions}

Let \(X\) be a nonzero normed space over \(\F=\mathbb R\) or \(\mathbb C\).
We say that \(X\) has \emph{countable capture} if for every countable
\(C\subset X\) there is \(u\in S_X\) such that
\[
        \dist(c,\F u)=\|c\|\qquad(c\in C).
\]

This is equivalent to the quotient-map and common-kernel formulations in the
earlier packet
\[
\texttt{2311.14257\_countable\_capture\_characterization}.
\]

We use the standard BCP form: \(X\) has BCP if and only if there are
countably many centres \(x_n\in X\setminus\{0\}\) such that
\[
        S_X\subseteq \bigcup_n B(x_n,\|x_n\|),
\]
where the balls are open.  Indeed, any ball not containing the origin has
radius strictly smaller than the norm of its centre, and enlarging its radius
to the centre norm gives the displayed form.  Conversely, \(0\notin
B(x_n,\|x_n\|)\) because the balls are open.

\section*{Main Theorem}

\begin{theorem}[exact BCP/capture equivalence]
\label{thm:main}
For every nonzero normed space \(X\), the following are equivalent:
\begin{enumerate}
    \item \(X\) fails BCP;
    \item \(X\) has countable capture.
\end{enumerate}
Equivalently, \(X\) has BCP if and only if there is a countable set
\(C\subset X\) such that for every \(u\in S_X\) there exists \(c\in C\) with
\[
        \dist(c,\F u)<\|c\|.
\]
\end{theorem}

\begin{corollary}[C*-algebras]
\label{cor:cstar}
For every \(C^*\)-algebra \(A\), the underlying Banach space of \(A\) has BCP
if and only if \(A\) does not have countable capture.  Hence \(A\) fails BCP
if and only if it has countable capture.
\end{corollary}

\section*{Idea of the Proof}

The one-way obstruction was already known in the run: if every countable
family of proposed centres admits a unit vector whose line preserves all
their distances, then that unit vector escapes every BCP ball.

The missing reverse implication is a scalar-rescaling trick.  If a fixed
countable set \(C\) blocks countable capture, then for every \(u\in S_X\)
some \(c\in C\) satisfies
\[
        \dist(c,\F u)<\|c\|.
\]
Thus \(c\) is closer than \(\|c\|\) to some scalar multiple \(\lambda u\).
The scalar \(\lambda\) is nonzero, and after dividing by it we get
\[
        \|u-\lambda^{-1}c\|<\|\lambda^{-1}c\|.
\]
So \(u\) lies in a standard BCP ball centred at a scalar multiple of \(c\).
A countable dense set of scalars keeps the family of such multiples
countable.

\section*{Proof}

\begin{proof}[Proof of Theorem~\ref{thm:main}]
First suppose \(X\) has countable capture.  Let \((x_n)\subset X\) be any
countable family of proposed BCP centres.  By countable capture, choose
\(u\in S_X\) such that
\[
        \dist(x_n,\F u)=\|x_n\|\qquad(n\ge1).
\]
Then, for every \(n\),
\[
        \|u-x_n\|\ge \dist(x_n,\F u)=\|x_n\|.
\]
Hence \(u\notin B(x_n,\|x_n\|)\) for every \(n\).  Since the centre family
was arbitrary, no countable BCP cover exists.  Thus \(X\) fails BCP.

Conversely, suppose \(X\) does not have countable capture.  Then there is a
countable set \(C\subset X\) such that for every \(u\in S_X\) there exists
\(c\in C\) with
\[
        \dist(c,\F u)<\|c\|.
\]
The element \(c\) cannot be zero.  Therefore for some scalar
\(\lambda\in\F\) we have
\[
        \|c-\lambda u\|<\|c\|.
\]
The scalar \(\lambda\) cannot be zero, since otherwise the inequality would
read \(\|c\|<\|c\|\).  Put \(\mu=\lambda^{-1}\).  Dividing by \(|\lambda|\)
gives
\[
        \|u-\mu c\|<\|\mu c\|.
\]

Let \(D\subset\F\setminus\{0\}\) be a countable dense set, for instance
\(\mathbb Q\setminus\{0\}\) in the real case and
\((\mathbb Q+i\mathbb Q)\setminus\{0\}\) in the complex case.  For fixed
\(u\) and \(c\), the set
\[
        \{\alpha\in\F:\|u-\alpha c\|<\|\alpha c\|\}
\]
is open in \(\F\) and contains \(\mu\).  Choose
\(q\in D\) in this open set.  Then
\[
        \|u-qc\|<\|qc\|.
\]
Thus every \(u\in S_X\) belongs to one of the balls
\[
        B(qc,\|qc\|),
        \qquad c\in C,\ q\in D.
\]
The set of such centres is countable, and all displayed open balls omit the
origin.  Hence \(X\) has BCP.

This proves both the equivalence \(X\) fails BCP if and only if \(X\) has
countable capture and the displayed equivalent positive formulation.
\end{proof}

\begin{proof}[Proof of Corollary~\ref{cor:cstar}]
Apply Theorem~\ref{thm:main} to the underlying complex Banach space of the
\(C^*\)-algebra \(A\).
\end{proof}

\section*{Consequences for the C*-Algebra Program}

The theorem proves the lane-19 conjecture
\[
        \text{\(C^*\)-algebra BCP failure}
        \quad\Longleftrightarrow\quad
        \text{countable capture}.
\]
In fact, no \(C^*\)-specific structure is needed for this equivalence.

The many previous negative packets now have a common interpretation: they
produce countable-capture witnesses.  The positive packets produce explicit
countable sets that block capture, often with uniform constants.  For example,
the countable product theorem gives UBCP by constructing a countable family
of centres \(c\) for which every unit vector satisfies a strict inequality
\(\|u-c\|<\|c\|\).

\section*{Limitations}

This packet settles the exact BCP/capture equivalence.  It does not by itself
settle the stronger quantitative question of UBCP.  A parallel quantitative
version would need uniform control of the gap between a covering radius and
the norm of the centre, plus boundedness of the radii.  The earlier
\(C^*\)-algebra positive packets provide UBCP in their model classes, but this
theorem does not prove that every \(C^*\)-algebra with BCP has UBCP.

\section*{Novelty and Search Bounds}

Local indexes were searched for arXiv:2311.14257, BCP, UBCP, countable
capture, quotient-line characterization, Birkhoff--James orthogonality, and
the earlier lane-19 capture packets.  The previous packet
\[
\texttt{2311.14257\_countable\_capture\_characterization}
\]
proved the quotient and dual forms of countable capture, while
\[
\texttt{2311.14257\_countable\_capture\_bcp\_obstruction}
\]
proved only the forward implication ``countable capture implies failure of
BCP.''  The reverse implication in this packet is the new scalar-rescaling
step.  Web searches during the surrounding June 17, 2026 pushes found the
source BCP literature and adjacent operator-space papers, but no prior record
of this exact capture/BCP equivalence.

\section*{Human Review Focus}

Reviewers should check:
\begin{enumerate}
    \item the negation of countable capture;
    \item the fact that \(\lambda\ne0\) in the strict distance inequality;
    \item the passage from \(\|c-\lambda u\|<\|c\|\) to
          \(\|u-\lambda^{-1}c\|<\|\lambda^{-1}c\|\);
    \item the use of a countable dense scalar set to keep the final BCP
          centres countable;
    \item the stated limitation that this is an exact BCP theorem, not a
          UBCP classification.
\end{enumerate}

\begin{thebibliography}{9}

\bibitem{HKL}
Jinghao Huang, Karimbergen Kudaybergenov, and Rui Liu,
\emph{The ball-covering property of von Neumann algebras and noncommutative
symmetric spaces},
arXiv:2311.14257, 2023.

\bibitem{lane19capture}
Run \texttt{fa\_banach\_001}, lane 19,
\emph{When the Countable-Capture BCP Obstruction Exists},
solution packet
\texttt{2311.14257\_countable\_capture\_characterization}, 2026.

\bibitem{lane19obstruction}
Run \texttt{fa\_banach\_001}, lane 19,
\emph{A Countable-Capture Obstruction for the Ball-Covering Property},
solution packet
\texttt{2311.14257\_countable\_capture\_bcp\_obstruction}, 2026.

\end{thebibliography}

\end{document}
